\section{Comparative Empirical Analysis: Prototyping vs. Traditional Techniques}
\label{sec:comparative}

This section delivers a head-to-head empirical evaluation addressing \textbf{RQ6.2} (Defect detection and disambiguation efficacy compared against textual specifications, JAD, and unstructured interviews).

\subsection{The Anchor Empirical Evidence: Rueda et al. (2020)}
The most rigorous empirical benchmark directly comparing prototyping against established elicitation methodologies is the family of experiments conducted by Rueda, Panach, and Distante \cite{rueda2020requirements}. Spanning four controlled replications with 167 software engineering students and practitioners grouped into 30 development teams, the study systematically evaluated three treatments:
\begin{enumerate}
    \item \textbf{Unstructured Interview (UI):} Traditional verbal questions without visual artifacts or structured group facilitation.
    \item \textbf{Joint Application Design (JAD):} Highly structured group elicitation workshops mediated by a trained facilitator with formal agendas.
    \item \textbf{Paper Prototyping (PP):} Collaborative exploration utilizing interactive paper interface mockups manipulated by stakeholders.
\end{enumerate}

\subsection{Statistical Findings on Missing and Unclear Requirements}
The empirical measurements from Rueda et al. \cite{rueda2020requirements} provide foundational insights into requirements quality metrics:

\begin{itemize}
    \item \textbf{Total Elicitation Volume (Discovery of Missing Requirements):} Paper Prototyping generated a statistically significant increase in the total volume of requirements captured compared to both JAD and Unstructured Interviews ($p = 0.02$). The tactile manipulation of mockups stimulated stakeholders to volunteer unprompted functional desires, successfully surfacing latent expectations that remained dormant during interview dialogues.
    \item \textbf{Specification Completeness ($p = 0.17$):} Despite the substantial increase in gross requirement volume, Paper Prototyping showed \textit{no statistically significant superiority} over JAD regarding specification completeness ($p = 0.17$). The empirical explanation lies in the selective visual bias: stakeholders thoroughly explored interactive screens, but frequently overlooked non-visual algorithmic constraints, background batch operations, and system-level exceptions.
    \item \textbf{Specification Correctness ($p = 0.59$):} Requirements correctness did not exhibit a statistically significant difference between PP, JAD, and UI ($p = 0.59$). The presence of a paper prototype does not inherently prevent stakeholders from misunderstanding underlying domain logic if the analyst does not challenge underlying assumptions.
    \item \textbf{Time Efficiency and Variance ($p = 0.08$):} Paper Prototyping required higher upfront preparation time to fabricate interface artifacts. However, during the elicitation session itself, PP demonstrated lower conversational variance and reduced the duration of contentious stakeholder debate compared to unstructured sessions.
\end{itemize}

\subsection{Methodological Validation and Confounding Variables}
In a subsequent replication, D{\'\i}az et al. \cite{diaz2021gender} expanded the experimental framework across multiple cohorts to investigate whether participant gender or interpersonal dynamics influenced elicitation outcomes across methods. Their findings demonstrated that Paper Prototyping established an egalitarian collaborative platform: unlike unstructured interviews, where vocal participants often monopolized conversational turn-taking, visual and tangible mockups democratized participation, ensuring that contributions remained invariant to gender distribution.

\subsection{Head-to-Head Comparative Synthesis}
Table~\ref{tab:comparative_matrix} presents a comparative synthesis of the primary elicitation and specification techniques based on the aggregated literature corpus.

\begin{table*}[t]
\centering
\caption{Comprehensive Empirical Comparison of Requirements Elicitation and Specification Techniques (RQ6.2)}
\label{tab:comparative_matrix}
\begin{tabularx}{\textwidth}{l X X X X}
\toprule
\textbf{Evaluation Dimension} & \textbf{Textual Specifications (SRS)} & \textbf{Unstructured Interviews} & \textbf{Joint Application Design (JAD)} & \textbf{Prototyping (Paper / Interactive)} \\
\midrule
\textbf{Missing FR Discovery} & Poor. Relies on abstract stakeholder imagination; high omission rate. & Moderate. Constrained by interviewer question prompts. & High. Structured group dynamics surface divergent business needs. & \textbf{Very High ($p = 0.02$).} Tangible navigation exposes missing workflows spontaneously \cite{rueda2020requirements}. \\
\addlinespace
\textbf{Ambiguity Resolution} & Poor. Natural language text hides semantic polysemy and vagueness \cite{cascon2025ambiguity}. & Low-to-Moderate. Verbal agreements frequently mask unaligned mental models \cite{ferrari2016ambiguity}. & High. Real-time stakeholder alignment resolves terminology conflicts. & \textbf{Very High.} Concrete UI layouts and state transitions eliminate lexical ambiguity. \\
\addlinespace
\textbf{NFR Coverage} & Moderate-to-High. Explicit standard NFR checklists are commonly employed. & Low. Stakeholders rarely articulate quality constraints spontaneously. & Moderate. Group debate identifies systemic security/reliability goals. & \textbf{Low (Passive) / High (LPP/Video).} Passive prototypes miss NFRs; LPP \cite{abad2018loud} and Video \cite{bano2023videos} rectify this. \\
\addlinespace
\textbf{Upfront Preparation Cost} & Minimal initial artifact preparation; high subsequent documentation time. & Minimal. Requires only basic topic outlines. & Moderate. Requires formal facilitation guides and schedule coordination. & \textbf{Moderate-to-High.} Requires artifact fabrication (paper cutouts or digital mockups). \\
\addlinespace
\textbf{Stakeholder Engagement} & Passive (``Sign-off on document'' syndrome). & Interpersonal dialogue; vulnerable to dominant personalities. & Collaborative structured workshop; high organizational commitment. & \textbf{Active \& Experiential.} Hands-on physical or digital manipulation \cite{diaz2021gender}. \\
\bottomrule
\end{tabularx}
\end{table*}

\subsection{Key Empirical Synthesis}
The empirical evidence decisively indicates that prototyping is superior to textual specifications and unstructured interviews for discovering missing functional requirements and resolving ambiguous interface behaviors. However, the data disproves the naive assumption that prototyping is an all-purpose silver bullet: without augmentation (e.g., Think-Aloud protocols or structured checklists), passive visual prototypes leave non-functional requirements and non-visual data integrity rules vulnerable to omission.
